% Generated draft: requires mathematical and source review.
\documentclass[11pt]{article}
\usepackage[margin=25mm]{geometry}
\usepackage{fontspec}
\setmainfont{Latin Modern Roman}
\usepackage{amsmath,amssymb,mathtools,unicode-math}
\setmathfont{Latin Modern Math}
\usepackage{enumitem,needspace,xurl,hyperref}
\setlength{\parindent}{0pt}
\setlength{\parskip}{5pt}
\setlength{\emergencystretch}{3em}
\newcommand{\N}{\mathbb{N}}
\newcommand{\Z}{\mathbb{Z}}
\newcommand{\Q}{\mathbb{Q}}
\newcommand{\R}{\mathbb{R}}
\newcommand{\C}{\mathbb{C}}
\newcommand{\F}{\mathbb{F}}
\newcommand{\A}{\mathbb{A}}
\newcommand{\PP}{\mathbb{P}}
\newcommand{\E}{\mathbb{E}}
\newcommand{\eps}{\varepsilon}
\newcommand{\dd}{\,\mathrm{d}}
\begin{document}
\section*{UnsolvedMath ID: 20000416}
\textbf{Source problem code:} catalogue entry
\subsection{Definitions and mathematical statement}
For a smooth projective complex curve \(C\), a Belyi map is a finite nonconstant morphism \(\beta:C\to\PP^1_{\C}\) whose branch values lie in \(\{0,1,\infty\}\). A Jacobian \(J\) of dimension \(g\) has complex multiplication (CM) when \(\operatorname{End}(J)\otimes\Q\) contains a commutative semisimple \(\Q\)-algebra of dimension \(2g\). A regular Belyi map is a Galois cover. These are the standard characteristic-zero conventions for the terms used in the source.

The supplied source is comment 4 of a problem attributed to Y. Zarhin. D. Ramakrishnan asks whether expressing ``this equation'' as a cover of \(\PP^1\) ramified only at \(0,1,\infty\) helps determine which of its Jacobians have CM, referring to work of Wolfart. The equation and preceding parts of the problem are absent from this record. Consequently the family and the intended meaning of reparametrization cannot be reconstructed from the original sentence alone.

The later report uses an unintroduced family \(C_{d,a}\). For integers \(d>7\), it states that the natural cyclic map has four branch values, so postcomposition by a M\"obius transformation of \(\PP^1\) cannot make that particular map Belyi. This restriction on the given map is distinct from existence of some Belyi map on the curve. Let \(C_d\) denote its asserted cyclic deck group of order \(d\) and \(N\) its normalizer in the automorphism group. In the family intended by the report, it claims
\[
 1\longrightarrow C_d\longrightarrow N\longrightarrow
 \operatorname{AffStab}(\{0,1,a\})\longrightarrow1,
\]
where \(\operatorname{AffStab}(S)\) consists of affine transformations \(x\mapsto ux+v\), \(u\ne0\), taking \(S\) to itself. Its candidate criterion asserts that the natural-normalizer route gives a regular abelian Belyi quotient precisely for
\[
 a\in\{-1,\tfrac12,2,\zeta_6,\zeta_6^{-1}\},\qquad
 \zeta_6=e^{2\pi i/6},
\]
equivalently the three harmonic values or the two roots of \(a^2-a+1=0\). The report states that these cases give degree-\(2d\) or degree-\(3d\) regular Belyi quotients with abelian deck group and recover those five known CM parameters through Wolfart's criterion. It claims four explicit quotient formulas but supplies none here.

The report explicitly restricts this to the natural-normalizer route, not a classification of all Belyi maps or all CM fibers. In its intended family, unrestricted Belyi-map existence for fixed \(d>7\) is reported equivalent merely to algebraicity of \(a\), or algebraicity modulo the anharmonic orbit (the orbit generated by \(a\mapsto1-a\) and \(a\mapsto1/a\)). It further reports that a corrected genus calculation and Moonen's classification identify \(d=9\) as a special genus-seven family with infinitely many additional CM fibers. Thus the five displayed parameters are not claimed to exhaust CM parameters. Because the equation, base-field conventions, allowable parameters (including the expected exclusions of collisions such as \(a=0,1\)), and full cited criteria are missing, these are recorded as supplied family-dependent claims, not a fully specified theorem about a newly invented \(C_{d,a}\).
\subsection{Short English statement}
The source asks whether a Belyi presentation helps detect CM Jacobians, but omits the equation being discussed. The supplied report distinguishes a restricted normalizer construction from arbitrary Belyi maps and from CM classification, including additional reported CM fibers when \(d=9\). The missing family equation prevents an independent precise family statement.
\subsection{Sources}
\begin{itemize}
\item Local UnsolvedMath registry, permanent ID \(20000416\), source code \texttt{catalogue entry}; AIM algebraic number theory corpus, \texttt{aim-algebraic-number-theory-notes.json}, notes[67]. Attribution: Ulam AI / UnsolvedMath Contributors, CC BY 4.0.
\item AIM workshop \emph{The Tate conjecture}, Problem 3 (Y. Zarhin), comment 4 attributed to D. Ramakrishnan: \url{https://aimath.org/WWN/tateconjecture/tateconjecture.pdf}. The comment mentions Wolfart. Original remarks and literature fields are empty; the verification text ends after ``The canonical record reads:''.
\item The research summary invokes Wolfart's criterion and Moonen's classification without full bibliographic references, and names \(C_{d,a}\) without giving its equation. No missing references or equation are invented.
\item Registry dated August 14, 2026: statement exact, status \texttt{partially\_solved}, classification \texttt{PARTIAL-PROGRESS}; attempt 1 labelled partial result, rigor proved, novelty confidence low. The normalizer criterion, quotient formulas, genus correction and CM assessment are machine-generated research aids, not peer reviewed, requiring independent verification before citation.
\end{itemize}
\end{document}
